\newpage
\begin{center}
  \textbf{\large АННОТАЦИЯ}
\end{center}

Работа посвящена моделированию динамики всплытия одиночного двумерного пузырька
методом физико-информированных нейронных сетей (PINN). В качестве источника
обучающих данных использованы результаты вычислительной гидродинамики (CFD),
полученные в пакете Basilisk для пяти значений начального радиуса пузырька.
Разработаны и исследованы два класса моделей: простая PINN, обученная на одном
значении радиуса пузырька, и параметрическая PINN, принимающая значение радиуса пузырька как дополнительный входной параметр и обобщающая решение на диапазон значений без переобучения.

Для параметрической модели проведено систематическое сравнение трёх архитектур PINN
(лёгкая, средняя, тяжёлая), исследовано влияние числа стадий обучения и проведено исследование влияния размера обучающей выборки на качество получаемого решения. Выявлено, что модель с параметрами (8$\times$350, $\tanh$, 2 стадии)
в тестах ($R\in\{0{,}20;\,0{,}30\}$) позволяет получить решение задачи с ошибкой $L_2$ порядка 6--11\% для скорости течения и менее 1{,}5\% для значения давления.
После однократного обучения предсказание 100 кадров занимает 1--7~с,
тогда как CFD-расчёт тех же кадров в Basilisk требует 11{,}7~с.
При этом время CFD-расчёта для каждого нового значения радиуса пузырька составляет несколько часов. Таким образом, параметризованная PINN-модель позволяет значительно ускорить процесс получения решения с приемлемой точностью.

\bigskip
\noindent\textbf{Ключевые слова:} физико-информированные нейронные сети,
многофазные течения, всплытие пузырька, CFD, суррогатная модель,
параметрическая нейронная сеть.

\onehalfspacing
\setcounter{page}{2}

\newpage
\renewcommand{\contentsname}{\centerline{\large СОДЕРЖАНИЕ}}
\tableofcontents

\newpage
\begin{center}
  \textbf{\large ВВЕДЕНИЕ}
\end{center}
\addcontentsline{toc}{chapter}{ВВЕДЕНИЕ}

Многофазные течения широко встречаются в химической промышленности, энергетике
и биотехнологиях: барботажные колонны, парогенераторы, ферментеры, флотационные
установки --- во всех этих системах ключевую роль играет динамика пузырьков.
Детальный расчёт таких течений традиционными методами вычислительной
гидродинамики требует значительных машинных ресурсов, что делает
параметрический анализ и оптимизацию в реальном времени практически недоступными.
В то же время современные производства требуют наличия оперативного контроля протекания процесса
и управление его ходом для обеспечения требуемого качества получаемого продукта и гарантирования бесперебойной работы технологических установок.

Физико-информированные нейронные сети (PINN) предлагают альтернативный подход к решению задач математической физики. Например, в области гидродинамики
нейронная сеть в ходе обучения учитывает в функции потерь уравнения Навье--Стокса и граничные условия для заданного процесса. В итоге физико-информированная нейронная сеть выполняет роль аппроксимации решения поставленной математической задачи. Это позволяет строить
непрерывные суррогатные модели, пригодные для быстрой оценки физических полей
при новых значениях параметров задачи.

Построение таких моделей часто связано с необходимостью исследования устойчивости модели, а также определения оптимальных гиперпараметров PINN, как с точки зрения точности получаемого решения, так и с точки зрения снижения вычислительной сложности решения.

\medskip
\noindent\textbf{Цели и задачи работы:}

В связи с вышесказанным целью настоящей работы является построение и оптимизация численной схемы решения задачи о всплытии пузырька на основе физико-информированной нейронной сети.

Для достижения поставленной цели требуется решить следующие задачи:

\begin{itemize}
  \item разработать архитектуру физико-информированной нейронной сети для решения двумерной задачи всплытия пузырька заданного радиуса в жидкости, позволяющую получить поля скорости, давления и индикатора фазы;
  \item выполнить параметризацию полученного решения для обобщения PINN-модели на различные
        значения начального радиуса пузырька;
  \item выполнить анализ влияния гиперпараметров модели на качество решения и время его получения и определить оптимальные характеристики архитектуры физико-информированной нейронной сети;
  \item оценить качество итоговой модели на тестовых радиусах, не вошедших в
        обучающую выборку, и сопоставить с результатами работы-прототипа на основе CFD-решателя.
\end{itemize}

\medskip
\noindent\textbf{Структура работы.}
Глава~1 содержит литературный обзор: обоснование актуальности задачи,
основные уравнения многофазной гидродинамики, классические численные методы
CFD и методологию физико-информированных нейронных сетей.
Глава~2 представляет теоретическое описание постановки задачи, а также
архитектуры простой и параметрической PINN.
Глава~3 описывает практическую часть: результаты всех экспериментов и
сравнительный анализ.
Глава~4 обсуждает физическую интерпретацию полученных результатов,
ограничения метода и направления дальнейшего развития.
В заключении подводятся итоги работы.
